\section*{Обозначения и сокращения}
\addcontentsline{toc}{section}{Обозначения и сокращения}
\label{sec:Abbreviations}

В настоящей работе применяются следующие обозначения и сокращения:

\begin{description}[leftmargin=!,labelwidth=2.6cm,font=\normalfont]
    \item[ВКР] --- выпускная квалификационная работа;
    \item[ИНС] --- искусственная нейронная сеть;
    \item[ЦОС] --- цифровая обработка сигналов;
    \item[ОСШ (SNR)] --- отношение сигнал/шум (signal-to-noise ratio);
    \item[STFT] --- кратковременное преобразование Фурье (short-time Fourier transform);
    \item[SSL] --- самоконтролируемое обучение (self-supervised learning);
    \item[ASR] --- автоматическое распознавание речи (automatic speech recognition);
    \item[TTS] --- синтез речи по тексту (text-to-speech);
    \item[SE] --- улучшение речи (speech enhancement);
    \item[SR] --- восстановление речи (speech restoration);
    \item[USM] --- Universal Speech Model, многоязычный SSL-энкодер Google;
    \item[GAN] --- генеративно-состязательная сеть (generative adversarial network);
    \item[CNN] --- свёрточная нейронная сеть (convolutional neural network);
    \item[FiLM] --- feature-wise linear modulation, послойная аффинная модуляция признаков;
    \item[RIR] --- импульсная характеристика помещения (room impulse response);
    \item[MOS] --- усреднённая субъективная оценка качества (mean opinion score);
    \item[PESQ] --- perceptual evaluation of speech quality;
    \item[STOI] --- short-time objective intelligibility;
    \item[SI-SDR] --- scale-invariant signal-to-distortion ratio;
    \item[DNSMOS] --- неинтрузивная нейросетевая оценка качества речи;
    \item[WER] --- доля ошибочно распознанных слов (word error rate);
    \item[RTF] --- коэффициент реального времени (real-time factor);
    \item[GPU] --- графический процессор (graphics processing unit).
\end{description}
